% Copyright (c) 2026 Eric van der Merwe
% SPDX-License-Identifier: MIT
%
% Study notes for Revelation. Same form as notes-Luke.tex:
%
% \Note chapter:verse {phrase from the verse}
% commentary...
%
\Note 4:1 {door was opened in heaven}
Intended to highlight that everything to follow is a demonstration/manifestation of forthcoming things.

\Note 4:1 {voice which I heard was as it were of a trumpet}
Any voice out of heaven symbolises Divine Truth or influx. John's mind was elevated and his understanding was opened to higher degrees of heaven. "Coming into the Spirit" entails entering a spiritual state in which heavenly things are made manifest. In the Old Testament, trumpets signfiy when assemblies and arrangements take place (think of ordering an army or signalling crowds). 

\Note 4:1 {things to come hereafter}
Revelation deals with nothing other than the last judgement and the end state of the church. This phrase just confirms that nothing else is dealt with.

\Note 4:2 {throne}
Thrones in the Bible always imply judgement, and the authority to judge.

\Note 4:2 {one sat on the throne}
This is clearly the Lord, as is made clear in the following verses.

\Note 4:3 {stone}
Stones correspond to truths in their ultimate forms. Precious stones correspond with truths transparent from their goods in ultimate form.

\Note 4:3 {jasper}
Jasper is a precious stone, that comes in red and white variants, often with inclusions of both variants at the same time. Here, it is used in reference to its white variant. White in general corresponds with truth.

\Note 4:3 {sardine}
Sardis is a precious stone, that comes in red and white variants, often with inclusions of both variants at the same time. Here it is used in reference to its red variant. Red in general corresponds with good.

\Note 4:3 {rainbow round about the throne, in sight like unto an emerald}
Rainbows appear throughout heaven, often in multi-colored and single-colored variants. This one is an emanation of the Lord's Divine Wisdom into heaven. In this instance it is green because of the natural heaven in which the Last Judgement occured. The Last Judgement was not on anyone in the heavens already, but rather on people who had accumulated in those heavens who were in natural truths only, or were externally good, but internally evil, behaving according to natural laws and restrictions instead of from an internal sense of good and truth. 

\Note 4:4 {four and twenty}
The numbers 12 and 24 denote a wholeness of something. For example, when referring to all 12 tribes of Israel, the Word is referencing all the goods and truths to do with the Church, not one thing in particular. 

\Note 4:4 {seats}
The word used here is also translatable as 'throne'. The 24 thrones refer to all things of judgement.

\Note 4:4 {}
The Lord/the elders do not perform the judgement involved here, the Word does. John 12:47-48 says, "I have come not to judge the world, but to save it. The Word I have spoken shall judge"

\Note 4:4 {white raiment}
The judgement involved would come from the truths of the Word.

\Note 4:4 {crowns of gold}
Crowns in the Word signify wisdom; gold crowns signify wisdom from love, which spawns all things of the Word and the Church anyway.

\Note 4:5 {}
Lightnings, thunderings, and voices were also present at Mt. Sinai when Moses received the Law, as well as other places in the Word. The Apostles James and John corresponded to charity, and works from charity. All perception is from these, hence they were called the "Sons of Thunder".

\Note 4:5 {lightnings}
Lightnings correspond to enlightenments, or first realizations of truth.

\Note 4:5 {thunderings}
Thunderings correspond to perceptions, or first understandings of truth.

\Note 4:5 {voices}
Voices correspond to instructions from the Word.

\Note 4:6 {sea}
A sea in the Word corresponds to Divine truth at the boundaries or fringes; angels in the most general truths from the letter of the Word are said to exist within an atmosphere like the sea. Elsewhere in the Word the Lord 'dried up the sea of Babel', which meant to extinguish even the ultimate truths from the Church of the time.

\Note 4:6 {crystal}
The quality of that sea is from heaven, and hence, from Divine Truth.

\Note 4:6 {four beasts}
The four beasts in the midst/around the throne correspond to the Word in its ultimates (the natural sense of the Word), and its guards. Cherubs and seraphs in the Word always represent the guards of the word. For example, the guard at the gates of Eden had a flaming sword that turned this way and that. This corresponded to Divine truth in ultimates that can be turned without corrupting the internals.

\Note 4:6 {eyes}
Eyes in reference to man correspond to the understanding; when referencing the Lord, they correspond to Divine Wisdom.

\Note 4:7 {}
All the animals mentioned here are also mentioned in Ezekiel in the same way.

\Note 4:7 {lion}
Lions correspond to the power of the Word, or the ability to influence.

\Note 4:7 {calf}
Calfs correspond to affections for knowing truths.

\Note 4:7 {man}
Men correspond to wisdom insofar as we are really men.

\Note 4:7 {eagle}
Eagles correspond to cognitions and their resulting understandings from Divine Truth.

\Note 4:8 {}
The number 6 is the product of 2 and 3. 2 corresponds to everything of good, while 3 corresponds to everything of truth. 6 corresponds to their totality.

\Note 4:8 {wings}
Wings correspond to powers or protections, depending on their use in context (e.g. an eagle is using its wings in a 'power' context to lift itself, while frequently the Word describes the Lord covering us with his wings, in a protections context). The lion's wings correspond to the Word's powers of fighting evils and untruths; the calf's wings correspond to the Word's power of affecting the mind; the man's wings correspond to the power of the Word in discerning the nature of God; the eagle's wings correspond to the Word's power in recognizing goods and truths.

\Note 4:8 {full of eyes within}
This corresponds to the Divine Wisdom within the natural sense of the Word.

\Note 4:8 {Holy, holy, holy}
The Word is continually teaching truths about the Lord alone. Triplication implies that all holiness lies in Him alone, debunking a trinity of three persons.

